% 10_riproducibilita.tex — Capitolo 10: Riproducibilità e ingegneria del progetto.
% Documento didattico "Intelligenza artificiale nel progetto supervised_telemedicina".

\chapter{Riproducibilità e ingegneria del progetto}\label{cap:riproducibilita}

In questo capitolo cambiamo prospettiva: non più l'IA in sé, ma
l'\emph{ingegneria} che la circonda. Un modello che non si può riprodurre
non è scienza, è un aneddoto. Il progetto \progetto{} è costruito attorno a
un principio semplice: \emph{stessi input, stessi output, sempre}. Vediamo
come questo principio si realizza in pratica: seed e split congelati, hash
di provenienza, report deterministico, test automatici e una dashboard che
riusa le stesse figure di questo documento.

\section{Perché la riproducibilità conta}\label{sec:perche-riproducibilita}

La riproducibilità è la base della fiducia in un progetto di machine
learning. Se due persone eseguono la stessa pipeline e ottengono numeri
diversi, non sappiamo quale dei due risultati è giusto — e non sappiamo
nemmeno \emph{perché} differiscono. In un progetto didattico questo è
ancora più importante: chi legge deve poter verificare ogni affermazione
di questo documento eseguendo i comandi e ottenendo gli stessi numeri.

L'analogia da programmazione è immediata: un \emph{build riproducibile}
compila lo stesso sorgente nello stesso binario, sempre. Un progetto di
machine learning riproducibile fa lo stesso con i dati: lo stesso codice e
lo stesso seed producono lo stesso dataset, lo stesso training e lo stesso
report. Senza questa proprietà, le metriche del capitolo~\ref{cap:valutazione}
sarebbero solo numeri senza contesto.

\section{Seed e split congelati}\label{sec:seed-split-congelati}

Il punto di partenza è il seed. In \file{src/telemedicina\_supervised/config.py}
il valore di default è \codice{SEED\_DEFAULT = 41}. Da questo seed base
derivano i seed dei tre blocchi del dataset (capitolo~\ref{cap:dati}),
registrati in \file{data/processed/metadati.json}:

\begin{table}[h]
\centering
\caption{Lo split congelato del dataset: seed derivati e dimensioni.}
\label{tab:split-riproducibilita}
\begin{tabular}{lrrr}
\toprule
\textbf{Blocco} & \textbf{Seed} & \textbf{Campioni} & \textbf{Scartati} \\
\midrule
train & 41 & 40\,000 & 52\,504 \\
validation & 42 & 10\,000 & 13\,192 \\
test & 43 & 20\,000 & 26\,341 \\
\bottomrule
\end{tabular}
\end{table}

I tre blocchi vengono generati \emph{una volta} e salvati su disco in
\file{data/processed/} come file \codice{.npy}. Il test è \emph{congelato}:
durante il tuning degli iperparametri (capitolo~\ref{cap:training}) si
consultano solo train e validation, e il test viene usato una sola volta,
alla fine, per la valutazione ufficiale. Se il test venisse ri-generato o
ri-usato, le metriche finali perderebbero significato: il modello avrebbe
``visto'' il giudice prima dell'esame.

\begin{esempio}{Perché i seed derivati}
I seed 41, 42 e 43 non sono tre numeri a caso: derivano dal seed base 41
con incrementi successivi. In questo modo basta cambiare un solo valore in
\file{config.py} per rigenerare l'intera pipeline in modo coerente, e il
legame tra i blocchi resta documentato nel codice e nei metadati.
\end{esempio}

\section{Hash di provenienza}\label{sec:hash-provenienza}

Congelare i dati non basta: bisogna sapere \emph{quale versione del codice}
li ha prodotti. \file{data/processed/metadati.json} registra gli hash
sha256 (primi 12 caratteri) dei tre file che determinano i dati:

\begin{itemize}
  \item \codice{hash\_safety\_rules}: \codice{c50c6c31739d} — la fonte
        unica delle soglie cliniche;
  \item \codice{hash\_synthetic\_generator}: \codice{5342681af107} — il
        generatore dei campioni;
  \item \codice{hash\_teacher\_rules}: \codice{f6968285b204} — il teacher
        che assegna le etichette.
\end{itemize}

Il meccanismo è semplice e potente: se una soglia clinica cambia in
\file{safety\_rules.py}, l'hash cambia e i dati generati con la vecchia
versione non sono più confrontabili con i nuovi. Il guard test
\file{tests/test\_single\_source.py} verifica inoltre che nessuna soglia
numerica sia duplicata altrove: il dataset, il teacher e il safety gate
dipendono tutti dallo stesso modulo (capitolo~\ref{cap:teacher}).

\section{Il report deterministico}\label{sec:report-deterministico}

Tutte le metriche di questo documento provengono da un unico file:
\file{data/models/report.json}. Il report è \emph{deterministico}: con lo
stesso seed e lo stesso codice, ogni esecuzione del training produce
esattamente gli stessi numeri — dalla loss di validazione
(\codice{loss\_val\_migliore = 0.05095}) alle metriche del test congelato
(accuracy $0.9851$, kappa $0.9775$), fino alla distribuzione delle classi
per blocco.

Questo è il risultato concreto della riproducibilità: i numeri citati nei
capitoli~\ref{cap:valutazione} e~\ref{cap:distillazione} non sono
un'istantanea fortunata, ma l'esito riproducibile di una pipeline. Chi
riesegue il training ottiene lo stesso \file{report.json}, e può
verificare ogni affermazione di questo documento.

\section{I test automatici}\label{sec:test-automatici}

La riproducibilità non riguarda solo i numeri: riguarda anche il
\emph{comportamento}. Il progetto ha 129 test automatici, eseguibili con
un solo comando:

\begin{lstlisting}[caption={Esecuzione della suite di test del progetto.}]
PYTHONPATH=src python -m unittest discover -s tests
\end{lstlisting}

La suite copre i moduli del runtime (MLP, scaler, safety rules, agente,
database), la pipeline di training, il contratto dei dati e la dashboard.
Due famiglie di test meritano attenzione:

\begin{itemize}
  \item \textbf{Guard test di dipendenza}: \file{tests/test\_dashboard.py}
        verifica che il runtime (\file{src/}, \file{training/},
        \file{main.py}) \emph{non} importi \codice{matplotlib} né i tool
        della dashboard. Il vincolo del progetto è esplicito: solo numpy e
        stdlib nel runtime, matplotlib solo in sviluppo.
  \item \textbf{Guard test di fonte unica}: \file{tests/test\_single\_source.py}
        verifica che le soglie cliniche vivano solo in
        \file{safety\_rules.py} e non siano duplicate altrove.
\end{itemize}

C'è anche una regressione sul codice legacy: i 14 test di
\file{archive/legacy\_qtable\_rule\_based/tests/} garantiscono che il
comportamento storico del sistema a regole non sia stato alterato durante
la migrazione al nuovo progetto.

\section{Architettura del codice}\label{sec:architettura-codice}

La riproducibilità è anche una questione di \emph{organizzazione}. Il
progetto separa nettamente le responsabilità in quattro aree:

\begin{itemize}
  \item \file{src/telemedicina\_supervised/}: il \textbf{runtime} — il
        codice che gira in produzione (safety rules, MLP, agente, database,
        servizi). Vincolo: solo numpy e stdlib.
  \item \file{training/}: la \textbf{pipeline} — generazione del dataset,
        teacher, training e metriche. Produce i file in \file{data/}.
  \item \file{tools/}: la \textbf{dashboard} — strumenti di sviluppo e
        visualizzazione, con matplotlib importato \emph{lazy} solo qui.
  \item \file{tests/}: la \textbf{verifica} — la suite di 129 test.
\end{itemize}

Questa separazione non è estetica: è ciò che rende possibile il guard test
sulle dipendenze. Se qualcuno importasse matplotlib nel runtime, il test
fallirebbe e la violazione del vincolo sarebbe visibile subito, non scoperta
per caso in produzione.

\section{La dashboard come strumento didattico}\label{sec:dashboard}

Chiudiamo con lo strumento che rende tutto visibile: la dashboard generata
da \file{tools/genera\_dashboard.py}. È un file HTML \emph{autonomo}
(\file{data/dashboard.html}) con sei viste:

\begin{enumerate}
  \item \textbf{Analisi Live}: la tabella delle analisi dal database;
  \item \textbf{Flusso decisionale}: il percorso di un singolo caso
        attraverso gate, MLP e fallback;
  \item \textbf{Teacher vs MLP}: metriche, matrice di confusione e regioni
        di decisione;
  \item \textbf{Training}: curve di loss e confronto della griglia;
  \item \textbf{Incertezza}: l'istogramma delle confidenze con la soglia;
  \item \textbf{Riproducibilità}: seed, split e hash da
        \file{metadati.json}.
\end{enumerate}

La dashboard riusa \emph{le stesse figure di questo documento}: la heatmap
delle regioni di decisione, l'istogramma delle confidenze, la matrice di
confusione e le curve di loss sono generate dagli stessi dati e con lo
stesso codice (\file{tools/\_figure.py}) che produce le immagini in
\file{figure/generate/}. Matplotlib è importato \emph{lazy} (solo se
serve) e vive esclusivamente in \file{tools/}: la dashboard funziona anche
senza, degradando le figure a segnaposto. È un esempio concreto del
principio di questo capitolo: la stessa pipeline, gli stessi dati, gli
stessi numeri — ovunque e per chiunque.